    \begin{question}{0.2}
        (20分) 已知某点处的应力张量为（单位：MPa）：
$$
\sigma_{ij} = \begin{bmatrix}
20 & 10 & 0 \\
10 & 30 & -15 \\
0 & -15 & -10
\end{bmatrix}
$$

(1) 应力张量分解
   将应力张量 \( \sigma_{ij} \) 分解为球张量 \(p\delta_{ij}\)（静水应力部分）和偏张量 \( s_{ij} \)（应力偏量部分）。即计算：
   $$
   p = \frac{1}{3}\sigma_{kk}, \quad s_{ij} = \sigma_{ij} - p\delta_{ij}
   $$

(2) 应力不变量与主应力
   a. 计算应力张量 \( \sigma_{ij} \) 的三个主应力不变量 \(I_1, I_2, I_3\)。
   b. 基于特征方程 \( \lambda^3 - I_1\lambda^2 + I_2\lambda - I_3 = 0 \)，求解其三个主应力 \(\sigma_1 \ge \sigma_2 \ge \sigma_3\)。
    \end{question}
    
    \begin{solution}
    (1) 应力张量分解
    \[
    p = \frac{1}{3}(20 + 30 + (-10)) = \frac{40}{3} \approx 13.33 \text{ MPa}
    \]
    \[
    s_{ij} = \sigma_{ij} - p\delta_{ij} = \begin{bmatrix}
    20 - \frac{40}{3} & 10 & 0 \\
    10 & 30 - \frac{40}{3} & -15 \\
    0 & -15 & -10 - \frac{40}{3}
    \end{bmatrix} \approx \begin{bmatrix}
    6.67 & 10 & 0 \\
    10 & 16.67 & -15 \\
    0 & -15 & -23.33
    \end{bmatrix} \text{ MPa}
    \]

(2)  应力不变量与主应力
    a.
    \[
    I_1 = \sigma_{kk} = 20 + 30 -10 = 40 \text{ MPa}
    \]
    \[
    I_2 = \begin{vmatrix} 30 & -15 \\ -15 & -10 \end{vmatrix} + \begin{vmatrix} -10 & 0 \\ 0 & 20 \end{vmatrix} + \begin{vmatrix} 20 & 10 \\ 10 & 30 \end{vmatrix} = (-300-225) + (-200) + (600-100) = -225 \text{ MPa}^2
    \]
    \[
    I_3 = \det(\sigma_{ij}) = 20(30\times(-10) - (-15)^2) - 10(10\times(-10) - (-15\times0)) + 0(...) = 20(-300-225) - 10(-100) = -10500 + 1000 = -9500 \text{ MPa}^3
    \]
    b. 特征方程为：\(\lambda^3 - 40\lambda^2 - 225\lambda + 9500 = 0\)。
    通过因式分解或数值求解可得三个主应力为：
    \[
    \sigma_1 \approx 43.8 \text{ MPa}, \quad \sigma_2 \approx 5.1 \text{ MPa}, \quad \sigma_3 \approx -8.9 \text{ MPa}
    \]
    \end{solution}
